\chapter{TÍCH HỢP VÀO LLM GATEWAY}
\label{chap:gateway-integration}

Dự án đích thực tế của chúng ta bao gồm giao diện Streamlit, dịch vụ web FastAPI, và sự hỗ trợ của hai nhà cung cấp mô hình dịch vụ (providers) lớn là DeepSeek và Gemini. Thay đổi mang tính bước ngoặt khi đưa bộ nhớ vào hệ thống không đơn thuần chỉ là thêm một lệnh import thư viện LangGraph. Đó là sự chuyển đổi toàn diện hợp đồng dịch vụ từ cơ chế ``gửi prompt đơn lẻ'' sang cơ chế ``gửi tin nhắn mới trong một phiên trò chuyện (conversation) được kiểm soát quyền sở hữu (owner)''.

\section{Hợp đồng API đề xuất cho cổng dịch vụ (LLM Gateway)}

Để đảm bảo tính độc lập và bảo mật, chúng ta đề xuất bộ khung tài nguyên API gồm các endpoint cơ bản sau:

\begin{lstlisting}[language=Python,caption={Tạo conversation và gửi message}]
# 1. End point tao phien tro chuyen moi
POST /conversations
{"user_id": "user-nhat"}

Response: 201 Created
{"thread_id": "c20f...", "user_id": "user-nhat"}

# 2. End point gui tin nhan moi vao phien tro chuyen
POST /conversations/c20f.../messages
{
    "user_id": "user-nhat",
    "message": "Ten toi la gi?"
}

Response: 200 OK
{"thread_id": "c20f...", "answer": "Ban ten la Nhat."}
\end{lstlisting}

\subsection{Luồng dữ liệu giữa Client và Server}
Giao diện người dùng (Streamlit frontend) chịu trách nhiệm lưu trữ định danh luồng hiện tại (\texttt{thread\_id}). Khi người dùng tạo cuộc trò chuyện mới, frontend sẽ gọi endpoint khởi tạo để nhận về một \texttt{thread\_id} dạng ngẫu nhiên (UUID). Ở các lượt chat tiếp theo, frontend chỉ gửi duy nhất tin nhắn mới nhất của người dùng kèm theo \texttt{thread\_id}; tuyệt đối không gửi kèm các tin nhắn cũ vì backend và checkpointer đã chịu trách nhiệm lưu giữ lịch sử.

\section{Xây dựng API hoàn chỉnh bằng FastAPI}

Dưới đây là mã nguồn dịch vụ FastAPI hoàn chỉnh tích hợp bộ nhớ hội thoại LangGraph:

\begingroup
\lstinputlisting[inputencoding=utf8/latin1,language=Python,caption={FastAPI conversation resource}]{code/fastapi_app.py}
\endgroup

\subsection{Phân tích chi tiết mã nguồn \texttt{fastapi\_app.py}}
\begin{itemize}
  \item \textbf{Cấu trúc kiểm tra quyền sở hữu \texttt{require\_owner}:} Hàm này lấy \texttt{thread\_id} và \texttt{user\_id} làm đối số. Nó truy vấn từ điển \texttt{thread\_owners} (đóng vai trò giả lập bảng CSDL hội thoại):
  \begin{itemize}
    \item Nếu \texttt{thread\_id} chưa từng được khởi tạo, hàm trả về lỗi HTTP 404 (Not Found).
    \item Nếu cuộc trò chuyện tồn tại nhưng thuộc quyền sở hữu của người dùng khác, hàm lập tức trả về lỗi HTTP 403 (Forbidden).
  \end{itemize}
  \item \textbf{Phương thức \texttt{create\_conversation}:} Khi người dùng gửi yêu cầu tạo luồng chat kèm theo định danh \texttt{user\_id}, hàm sinh một mã UUID ngẫu nhiên làm \texttt{thread\_id}, đăng ký quyền sở hữu vào từ điển \texttt{thread\_owners} và trả lại thông tin cho client.
  \item \textbf{Phương thức gửi tin nhắn \texttt{send\_message}:} Hàm thực hiện kiểm tra quyền sở hữu trước thông qua \texttt{require\_owner}, sau đó gọi phương thức \texttt{chat\_service.ask} để thực thi đồ thị LangGraph và nhận về câu trả lời.
  \item \textbf{Phương thức xóa hội thoại \texttt{delete\_conversation}:} Giải phóng toàn bộ dữ liệu checkpoint ngắn hạn trong checkpointer thông qua \texttt{chat\_service.delete\_thread} và xóa quyền sở hữu trong từ điển \texttt{thread\_owners}.
\end{itemize}

\section{Adapter tích hợp mô hình thực tế (DeepSeek và Gemini)}

Để ứng dụng dễ dàng chuyển đổi qua lại giữa hai nhà cung cấp dịch vụ DeepSeek và Gemini mà không làm xáo trộn cấu trúc đồ thị, chúng ta sử dụng các lớp bọc (wrappers) chuẩn hóa tin nhắn của LangChain:

\begingroup
\lstinputlisting[inputencoding=utf8/latin1,language=Python,caption={Adapter DeepSeek và Gemini}]{code/real_model_adapter.py}
\endgroup

\subsection{Phân tích chi tiết mã nguồn \texttt{real\_model\_adapter.py}}
\begin{itemize}
  \item \textbf{Bọc mô hình DeepSeek:} Sử dụng lớp \texttt{ChatOpenAI} từ gói \texttt{langchain\_openai}. Nhờ tính tương thích cao của API DeepSeek với chuẩn OpenAI, chúng ta chỉ cần chuyển đổi địa chỉ kết nối \texttt{base\_url} thành \texttt{https://api.deepseek.com} và truyền API key tương ứng.
  \item \textbf{Bọc mô hình Gemini:} Sử dụng lớp \texttt{ChatGoogleGenerativeAI} từ gói \texttt{langchain\_google\_genai} và truyền mã khóa API Google.
\end{itemize}

Khi có các adapter trên, Node gọi mô hình trong đồ thị sẽ cực kỳ ngắn gọn và độc lập với nhà cung cấp:

\begin{lstlisting}[language=Python,caption={Thay model quyết định bằng chat model thật}]
async def call_model(state: MessagesState):
    # Cut old messages based on token budget
    model_messages = context_for_model(
        state["messages"],
        max_tokens=8_000,
    )
    # Call model asynchronously via ainvoke
    response = await model.ainvoke(model_messages)
    return {"messages": [response]}
\end{lstlisting}

\begin{figure}[H]
\centering
\begin{tikzpicture}[
 box/.style={draw=aiBlue,rounded corners,fill=aiSoftBlue,minimum width=3.05cm,minimum height=1cm,align=center},
 arr/.style={-{Stealth[length=2mm]},thick,aiNavy}]
 \node[box] (ui) {Streamlit\\message + thread};
 \node[box,right=0.8cm of ui] (api) {FastAPI\\ownership};
 \node[box,right=0.8cm of api] (graph) {LangGraph\\state + Store};
 \node[box,right=0.8cm of graph] (llm) {DeepSeek\\hoặc Gemini};
 \draw[arr] (ui)--(api); \draw[arr] (api)--(graph); \draw[arr] (graph)--(llm);
 \draw[arr] (llm) to[bend right=28] (ui);
\end{tikzpicture}
\caption{Ranh giới thành phần trong LLM Gateway có memory.}
\label{fig:gateway-architecture}
\figsource{Tác giả tự dựng.}
\end{figure}

\section{Lộ trình di cư mã nguồn (Migration Slices)}

Để tích hợp hệ thống bộ nhớ vào một dự án Gateway lớn đang chạy ổn định một cách an toàn nhất, chúng ta cần chia nhỏ quy trình thành các lát cắt nhỏ (thin slices):

\begin{enumerate}
  \item \textbf{Lát 1:} Giữ nguyên mô hình giả (rule-based mock model); tích hợp cổng FastAPI để kiểm thử cô lập tính cách ly và bảo mật giữa các luồng.
  \item \textbf{Lát 2:} Thay thế mô hình giả bằng mô hình thực thông qua adapter bọc; chưa tích hợp cơ chế tóm tắt (summary) hay cắt bớt (trimming).
  \item \textbf{Lát 3:} Tích hợp hàm cắt bớt tin nhắn \texttt{trim\_messages} và các metric đo đạc token đầu vào/đầu ra.
  \item \textbf{Lát 4:} Thay thế lưu trữ tạm thời trong RAM bằng lưu trữ SQLite hoặc PostgreSQL thực tế; kiểm thử kịch bản tắt/mở server.
  \item \textbf{Lát 5:} Bổ sung tính năng lưu trữ tùy chọn cá nhân dài hạn (Store) khi người dùng chủ động cho phép.
\end{enumerate}

\begin{aiwarning}[Tránh tích hợp mọi tính năng cùng một lúc]
Nếu bạn đồng thời thay đổi cấu hình cổng API, thay thế adapter gọi mô hình, cấu hình thêm database lưu trữ checkpoint, và bật tính năng tóm tắt cùng một lúc, khi xảy ra lỗi (như chatbot trả lời sai hoặc mất lịch sử), bạn sẽ hoàn toàn không thể định vị được nguyên nhân nằm ở tầng ngữ cảnh (prompt), tầng lưu trữ (persistence), hay lỗi logic của mô hình. Mỗi bước di cư cần được đi kèm với các bài viết kiểm thử tự động tương ứng.
\end{aiwarning}

\begin{aicheckpoint}[Lab dự án 30 phút]
Khởi chạy dịch vụ FastAPI local bằng lệnh: \texttt{uv run uvicorn fastapi\_app:app --app-dir code --reload}. Sử dụng tài liệu tương tác Swagger API để:
\begin{enumerate}
  \item Khởi tạo hai cuộc trò chuyện độc lập của người dùng A và người dùng B.
  \item Thực hiện gửi tin nhắn khai báo tên khác nhau vào mỗi cuộc trò chuyện.
  \item Thử nghiệm dùng tài khoản của người dùng B để gửi yêu cầu đọc lịch sử tin nhắn của người dùng A. Kết quả mong muốn: Mã trạng thái lỗi bảo mật HTTP 403 trả về thành công trước khi hệ thống thực hiện đọc checkpoint.
\end{enumerate}
\end{aicheckpoint}
